\clearpage
\section{Twelve discriminating predictions and research gates}
\label{sec:predictions}
\subsection{Readiness, endpoints, and interpretation}
The following cards were introduced as unexecuted proposals in version 2. They remain the target definitions; version 3 executes feasible subsets in Section~\ref{sec:local-tests}, with exact coverage and remaining gaps reported there. Version 4 adds the targeted revision evaluations in Section~\ref{sec:revisions}; it preserves the original target definitions and historical outcomes. They are not public registrations or claims that every hypothesis will succeed. Z1--Z4 are engineering hypotheses; F1 is an exact negative consequence of a stipulated model; F2--F4 are conditional physical targets; T1--T2 are computational predictions; T3 is a consistency gate with no assumed successful outcome; T4 is a candidate empirical target requiring further matching. Their common form is an explicit observation, a comparator, and a result that would count against the claim.

Version 2 prioritized Z1, a prospective memory test after completing the isolated Z3 package, and an independent numerical T1 route. Version 3 now executes the bounded Z3 comparison and selected T1 routes; the full Z1 factorial and broader independent replication remain open. The most discriminating laboratory candidates are the joint separation--duration--state dependence in F2 and the shared-parameter relation in F3, but both remain conditional on an operational apparatus and independently constrained bath. A bare damping signal would be weak evidence because ordinary decoherence also damps visibility. T3 should precede claims that the spectral branch describes a physical vacuum; T4 requires a credible spectrum and threshold map first.

Each superiority claim should specify a scientifically meaningful effect $\delta$ before data collection and report its interval. For a paired loss difference $\Delta$ where smaller is better, a possible decision rule is an upper confidence bound below $-\delta$. A separate noninferiority guardrail requires its upper bound below a predeclared harm margin. An interval crossing the margin is inconclusive, not proof of equivalence or mechanism absence. For a one-sided normal approximation, a planning quantity is
\begin{equation}
 n_{\rm independent}\simeq
 \frac{(z_{1-\alpha}+z_{1-\beta})^2\sigma_d^2}{(\Delta_1-\Delta_0)^2},
\end{equation}
where $\Delta_0=-\delta$ is the superiority boundary, $\Delta_1<\Delta_0$ is the assumed true effect, and $\sigma_d^2$ is the independent-cluster paired-difference variance. It is not a numerical sample-size prescription for the four-record pilot. The required variance and meaningful margins need external justification or separate pilot data. Source clusters, multiple endpoints, repeated looks, and selection of the best baseline must enter the final design. Choose one confirmatory primary contrast per study or predeclare familywise control; retain exploratory analyses as exploratory.

\subsection{ZoraASI engineering cards}
\paragraph{Z1: Explicit evidence relations reduce avoidable abstention.}
\textbf{Prediction and endpoint:} the revised scope rules and relation schema reduce unnecessary abstention on answerable unseen claims without exceeding a prespecified increase in incorrect definite judgments. \textbf{Comparison:} frozen original versus revised single-pass systems, followed by a rules-by-schema factorial ablation. \textbf{Controls:} identical evidence and model digest, held-out records and dates, independent scope/truth labels, blinded adjudication, clustered inference, and failures included in an overall endpoint. \textbf{Decision:} support requires the abstention reduction and the false-certainty guardrail together. A coverage increase alone fails the claim. The existing twelve-case result selects a candidate; it does not estimate a transportable effect. \textbf{Readiness:} protocol specification and independent labels required; no novel physics required.

\paragraph{Z2: Blinded challenge helps through error diversity.}
\textbf{Prediction and endpoint:} hiding the verifier's verdict from the challenger lowers selective error at a prespecified common coverage. \textbf{Comparison:} blind versus informed challenger within the same three-role architecture, plus a single-pass baseline matched on actual resource budget. \textbf{Controls:} identical source families, randomized arm order, separate measurements of joint verifier/challenger mistakes, truncation, tokens, and latency. \textbf{Decision:} no error reduction at matched coverage/cost, or a benefit explained entirely by refusing more cases, fails the proposed advantage. Equal total output allowance does not by itself equalize repeated input cost or truncation risk. \textbf{Readiness:} engineering experiment; the saved pilot currently gives no three-role accuracy gain.

\paragraph{Z3: Adaptive forgetting improves the recurrence--change tradeoff.}
\textbf{Prediction and endpoint:} a causal mixture of forgetting rates improves equally weighted held-out scenario-average Brier loss relative to the best fixed rate selected on the same validation mixture, subject to worst-scenario and stationary-regime harm limits. \textbf{Comparison:} selected fixed rate, adaptive mixture, uniform mixture, no memory, and TTL. \textbf{Controls:} common paths, disjoint seeds, prequential scoring, identical shadow base/archive for weighting ablations, and separate replacement-policy factors. \textbf{Decision:} lack of improvement over the strongest selected baseline or violation of the guardrail fails the bounded claim. Phase 2.9's short-dwell counterexample rules out unconditional benefit. \textbf{Readiness:} the original Phase 2.10 bundle lacks a dependency and a complete test package. The isolated version 3 experiment supplies these locally and evaluates a specific frozen comparison; broader robustness remains open.

\paragraph{Z4: Provenance awareness resists duplicated and stale evidence.}
\textbf{Prediction and endpoint:} source-ancestry and time metadata reduce increases in false definite judgments under duplication or obsolete-report perturbations while preserving appropriate use of genuinely new evidence. \textbf{Comparison:} metadata present versus ablated with paired, truth-preserving perturbations: duplicates, independent corroboration, superseded reports, unavailable current records, and authentic but irrelevant quotations. \textbf{Controls:} equalized context/source budgets, blinded semantic labels, fixed underlying propositions, and record-level clustering. \textbf{Decision:} unwarranted certainty rises materially under duplicate/stale evidence, or independent corroboration cannot be distinguished from copying. Numerical confidence scoring is a new endpoint if added; it cannot be inferred from the existing verdict labels. \textbf{Readiness:} controlled evidence benchmark.

\subsection{Conditional physical cards}
\paragraph{F1: Fixed-count population bits are insensitive to pure dephasing.}
\textbf{Prediction:} under Proposition~\ref{prop:bits}, bit-value distributions and their likelihood have zero dependence on the dephasing rate. \textbf{Comparison:} retain path-basis bits versus recombined, phase-resolved interference data from the same stipulated dynamics. \textbf{Observable:} likelihood derivative or estimated rate information, not a post hoc bit imbalance. \textbf{Controls:} known preparation, readout basis, independence, fixed count, and timing exclusion; the interference model must retain nuisance contrast and phase. \textbf{Failure interpretation:} parameter dependence in the exact stipulated bit model rejects the implementation or an assumption. A laboratory bias suggests preparation, readout, population transfer, or another mechanism; it is not distinctive support for rate-only physics. \textbf{Readiness:} analytic and software test available now; it does not require a claimed field detection.

\paragraph{F2: Quadratic-portal decoherence has a constrained spatial shape.}
\textbf{Prediction:} the massive thermal weak-coupling point-probe model gives zero rate at coincident paths, quadratic small-separation scaling, a finite saturation value, and the stated low-temperature activation. \textbf{Observable:} $-\log(V/V_{\rm ref})/\tau$ over several separations, durations, and independently characterized bath states; a ratio at two nonzero rates removes the shared coupling. \textbf{Comparison:} the full kernel versus conventional noise mechanisms and a generic fitted damping curve, with identical nuisance treatment. \textbf{Decision:} a jointly incompatible shape rejects the frozen bath/portal/apparatus model after quantified systematics. Agreement does not uniquely identify a scalar. No signal cannot reject arbitrarily weak coupling or absent bath population. \textbf{Readiness:} conditional numerical curve supplied here; finite trajectories, composite form factors, bath preparation, and power remain open.

\paragraph{F3: One portal parameter must fit both width and visibility.}
\textbf{Prediction:} Eq.~\eqref{eq:cross} holds when the same scalar, normalization, mass, and portal account for Higgs decay and the prepared probe environment. \textbf{Observable:} a joint likelihood for the width and visibility kernel with independently constrained bath occupation and nuclear inputs. \textbf{Comparison:} a single shared coupling versus unrelated effective amplitudes, including ordinary decoherence. \textbf{Decision:} no allowed shared parameter region rejects the specified common explanation, conditional on nuisance and model adequacy. A width upper limit by itself supplies no lower bound on visibility loss. \textbf{Controls:} invisibility assumptions, thresholds, data overlap, total-width conventions, and absence of double counting. \textbf{Readiness:} algebraic bridge available; no joint empirical fit or new exclusion is reported.

\paragraph{F4: Calibrated rate modulation obeys a composition constraint.}
\textbf{Prediction:} if two calibrated source contributions add in $\mathcal E$, the exponential response satisfies
\begin{equation}
 \log\!\left[\frac{\gamma(\mathcal E_1+\mathcal E_2)\gamma(0)}
 {\gamma(\mathcal E_1)\gamma(\mathcal E_2)}\right]=0.
\end{equation}
\textbf{Observable:} a source-factorial rate contrast with positive rates and fixed channel/readout. \textbf{Comparison:} exponential, linear-response, saturating, and ordinary environmental models over a predeclared range. \textbf{Decision:} a contrast incompatible with zero rejects additive exponential response under the calibrated source assumptions. Small-signal agreement is weak discrimination because many models share a linear Taylor term. \textbf{Controls:} independent field/source calibration, baseline drift, nonlinear source transfer, nuisance rates, and held-out amplitude levels. \textbf{Readiness:} source and apparatus objects remain open; cognitive or ethical labels are not calibrated physical sources.

\subsection{Spectral and consistency cards}
\paragraph{T1: Exact finite-radius kernels split the local mixed quartics.}
\textbf{Prediction:} for the fixed supplied operator, profile, radius, and scalar convention, $C_0+C_1=2A_4$ while $C_1/C_0\simeq1.59179265$, not one. \textbf{Comparison:} independent block-exponential and Cauchy-contour implementations, increased angular cutoffs, and the local approximation. \textbf{Observable:} coefficient differences and cutoff convergence. \textbf{Proposed tolerances:} absolute coefficient agreement $10^{-9}$ and sum-identity residual $10^{-10}$, to be frozen alongside precision and cutoff rules. These tolerances are new protocol choices, not observed errors. \textbf{Decision:} persistent disagreement after controlled convergence rejects the stated calculation; recovery of unequal coefficients does not reject it. \textbf{Readiness:} independent computational replication can proceed from the saved package. Physical masses require further normalization and vacuum matching.

\paragraph{T2: The protected photon direction survives the scalar response.}
\textbf{Prediction:} neutral alignment and the stipulated charge imply $M^2N=0$, while the kinetic residue follows the quoted quadratic response. \textbf{Comparison:} an independent full external-Clifford finite-momentum implementation versus the saved reduction. \textbf{Observable:} a scale-normalized mass-null residual, the scalar--gauge Ward relation, and a remainder scaling as $O(h^4)$ after subtracting the quadratic term under uniform field rescaling. \textbf{Decision:} a residual surviving precision/cutoff control or a remainder inconsistent with the $O(h^4)$ bound rejects the specified calculation. Roundoff near zero must be controlled. \textbf{Readiness:} computational test; photon masslessness is common to other gauge theories and is not distinctive empirical confirmation of this construction.

\paragraph{T3: A complete benchmark must pass coupled consistency.}
\textbf{Gate, not an assumed positive prediction:} in one frozen scheme solve joint radius/scalar stationarity and flatness, require positive physical kinetic eigenvalues and an acceptable normalized Hessian after removing gauge directions, then compute and cancel the appropriate higher-dimensional anomalies with stipulated fields/counterterms. \textbf{Comparison:} exact treatment versus controlled truncations inside their own domains. \textbf{Decision:} no admissible solution in a rigorously covered domain, a forbidden physical instability, or an uncancelled irreducible anomaly rejects that benchmark. An unsuccessful bounded numerical search is only a bounded search result. \textbf{Controls:} fixed allowed parameters and counterterms; no unregistered rescue term. \textbf{Readiness:} unresolved theoretical work. Success would establish conditional consistency, not observed unification.

\paragraph{T4: A fixed light spectrum predicts relative running slopes.}
\textbf{Prediction:} in a predeclared threshold-free interval with the stipulated active matter, Eq.~\eqref{eq:running} fixes the relative inverse-coupling slope. \textbf{Observable:} multi-scale scattering observables translated using a single matching scheme and covariance model. \textbf{Comparison:} the frozen joint slope vector versus alternatives with different independently specified spectra; constant matching offsets may float. \textbf{Decision:} incompatible slopes reject that spectrum/completion at the declared perturbative accuracy. Choosing thresholds from the same test data destroys the claimed out-of-sample prediction. \textbf{Controls:} physical masses, hypercharge normalization, thresholds, higher-order errors, and correlated datasets. \textbf{Readiness:} phenomenology and spectrum audit required, especially for the four chiral generations. No measured running fit is performed here.

\subsection{Explicit ledger of speculative connections}
\begin{longtable}{@{}>{\raggedright\arraybackslash}p{.25\linewidth}>{\raggedright\arraybackslash}p{.34\linewidth}>{\raggedright\arraybackslash}p{.34\linewidth}@{}}
\toprule
Speculative bridge & What the record supports & What would be needed\\\midrule
\endfirsthead
\toprule Speculative bridge & What the record supports & What would be needed\\\midrule\endhead
\bottomrule\endfoot
Functional self-monitoring implies experience & Forecasts and internal-state variables can be operationally scored. & An independent theory and validated discriminating measurements of experience; performance labels alone are insufficient.\\
Cognitive valence or coherence is $\mathcal E$ or $\varphi$ & Software and physiological quantities have domain-specific definitions. & Units, an observable field estimator, causal source manipulation, coupling action, controls, and replicated data.\\
Ethical intention sources a physical field & Normative constraints can govern choices and experiments. & A noncircular physical source definition and falsifiable controlled effect; a moral score is not a source term.\\
Narrative branch selection is quantum outcome control & The source explicitly uses metaphor/toy descriptions; the inspected rate branch preserves population probabilities. & A different explicit dynamics, consistency checks, and distinctive observations. It cannot be inferred from a fictional scene.\\
Machine memory resembles biological inheritance & Memory policies admit causal software ablations. & A biological mechanism, genotype/phenotype measurements, and independent evidence for any inherited architecture.\\
One unified consciousness--matter theory follows from common mathematics & Several branches use constraints, noise, feedback, and symmetry. & A single action and consistent limits, matched scales/parameters, identified observables, and successful independent tests. Shared vocabulary does not supply that bridge.\\
\end{longtable}

Keeping these entries explicit makes exploration possible without turning an analogy into a result. A speculative bridge can become a conditional model when its variables and equations are defined, and an empirical claim only when a suitable observation is supplied. It need not be deleted from the research agenda; it must retain its current evidential status.
